\subsection{大非川以后：边地并不等待中原改号}
大非川的失利在670年发生时，唐廷面对的并不是一条抽象的“西北边线”，而是青海高原、河西走廊、吐谷浑故地与通向安西的多组道路。战败的日期和唐方叙述可由\eventref{ST-670-DAFEICHUAN}{大非川之战}追索；兵数、将领动机和吐蕃内部的完整安排却不能由后来的汉文史书确定。更稳妥的判断是，距离、海拔、粮秣与对手的动员使唐军的选择受限，而这种限制会一直进入武周时期的用兵与人事。

692年恢复安西四镇，见\eventref{ST-692-ANXI-RECOVERY}{恢复安西四镇}，常被概括为“收复西域”。这个说法掩盖了绿洲、军镇和路程各不相同。收回名义上的据点需要军力，维持守军又要依靠仓储、马匹、地方征发与中转站；吐鲁番文书所保存的若干行政痕迹，正说明文书和人员必须在地方层面重新接合。它们并不能证明所有绿洲按同一制度运作，更不允许把城镇的短期控制画成没有缝隙的国界。

695年契丹首领李尽忠、孙万荣起事，见\eventref{ST-695-KHITAN-REBELLION}{契丹起事}，把河北与东北的道路、人事和边军问题带入中原政治。消息到达洛阳时，朝廷可以任命将领、下诏征兵；它仍要依靠沿路州县、军镇与临时合作者把命令变成行动。正史偏好记住朝廷如何惩奖主将，较少留下被征发者、向导或边地家庭的完整声音。承认这个空白，不是把边地写成被动背景，而是拒绝用中央的词汇替它们安排一致的立场。

\subsection{周的国号、寺院的语言与行政的连续}
690年武后改唐为周，见\eventref{ST-690-ZHOU-FOUNDATION}{武周建国}。改号经过诏令、官员请愿、宗室处置与符瑞文字，因而既是宫廷权力的重新布置，也是文书、年号和官衔必须被地方接受的行政操作。佛教文本在这一政治语言中取得位置，并不使所有寺院都成为同一意志的传声处；造像、写经和功德题记可让一些施主与僧侣露面，却只能说明特定地点、特定时段的实践。

705年复唐，见\eventref{ST-705-TANG-RESTORATION}{恢复唐号}，也不能被读作一夜之间把周的行政活动从社会中抹去。墓志中的年号、官衔与家族叙述有时能显示人们如何在转换中安置自身，但墓志首先是家族希望保存的公共记忆。它们可以校正某个人的任官或死亡日期，不能替我们测量整座州城的认同。复唐的意义在于宫廷控制和合法性语言改变；它的地方速度仍取决于传递诏令的人、继续办事的吏员以及各地可供调动的资源。

712年玄宗即位、713年平太平公主，见\eventref{ST-712-XUANZONG-ACCESSION}{玄宗即位}与\eventref{ST-713-TAIPING}{平太平公主}。两件事都发生在宫廷近处，却不能只用宗室性格解释。禁军、门禁、官员网络与谁掌握消息，使继承从家族关系变成可实行或可失败的政治行动。继承争夺结束后，朝廷才较能把注意力投向户籍、仓储、运河和边地；但这些事务并不会因宫廷安静而自行顺畅。

\subsection{开元的粮账不是盛世的同义词}
宇文融在721年的括户，见\eventref{ST-721-YUWEN-RONG}{宇文融括户}，把国家财政问题带到村落、土地与流动人口的交界。中枢希望辨认漏籍、占田与应税资源，地方官却必须在旧帐、亲属关系、逃亡与豪强的实际力量中办事。后世史书记录的增户、增田和成效带有奏报与编纂的口径，不能直接换成全国人口或农民生活的精确统计。敦煌、吐鲁番的文书能让一些登记和赋役的程序可见，恰恰因为它们只属于有限地点，才提醒我们地区差异不能被总数吞没。

裴耀卿在733年经营转运，见\eventref{ST-733-PEI-YAOQING}{裴耀卿转运}，接着处理另一个问题：被辨认出来的资源如何在水位、季节、仓场与船户的约束下送到需要它的地方。江淮水路不只是中央财政的背景，它也是沿线人群看见征调、交易、雇役和风险的空间。制度文字能够告诉我们某项安排被命名为何物；它不能证明每一船粮都按规划到达，也不能把一处仓场的遗存外推为全国的富足。

751年怛罗斯战事与安禄山兼领三镇，让这种表面上的整合显出另一面。西向的军事行动受制于更远的补给和地方联盟，东北的将领则把多个边镇的兵员、将校与财政机会集中到同一指挥体系。前者见\eventref{ST-751-TALAS}{怛罗斯战事}，后者见\eventref{ST-751-AN-LUSHAN-COMMANDS}{安禄山兼领三镇}。二者没有一条自动相连的因果线；它们共同提示的是，开元天宝时期的国家能力越能远距调度，失去平衡时能被集中起来的风险也越大。
